\subsection{Capacity of the network}
A first validation test of the implementation is to reproduce the evolution of the reconstruction error shown in \Cref{fig:phase_transition}. To this end, we choose the simplest case that is to initialize the system directly on one of the stored patterns and checking whether the network is able to retain the memory of that pattern throughout its evolution. This avoids having to fix, as an additional parameter, the initial distance between the system and the stored pattern. The model is initialized on a $32 \times 32$ fully connected square lattice. Each replica possess its own independend stored patterns. We fix the temperature at $\beta=5.0$ in order to approach the $T\to 0$ limit case and we let the system evolve for $2^{14}$ sweeps, averaging on the second half of the evolution. We show in \Cref{fig:reconstruction_error} the obtained reconstruction error depending on $\alpha$.
\nolinenumbers
\begin{figure}[htbp]
\centering
\includegraphics[width=\columnwidth]{reconstruction_error_vs_alpha_multiN.pdf}
\caption{Evolution of the reconstruction error with the storage load for several lattice sizes. For this diagram, the configuration is initialized on one of the stored patterns.}
\label{fig:reconstruction_error}
\end{figure}
\linenumbers
We observe the expected evolution, with a null reconstruction error for the lowest $\alpha$ and a burst increase around the critical value before a slowly increasing plateau, where the system fails to retain the memory of the initially stored pattern. The increase is steeper for the largest systems, the finite size effects being less and less significant. In this regime, the plateau value tends to increase with the lattice size, approaching the expected reconstruction error of 0.5, obtained when the global orientations of the neurons are totally uncorrelated with that of the stored patterns. The first explanation for this observation is that initializing the network on the exact pattern favors better retrieval. This explanation is however not robust to the increase of the plateau with the storage load. The second is that finite size effects favor a better random alignment, as the reconstruction error is computed based on the best overlap, and that smaller systems favors extreme value, the noise on each overlap being of the order of $\sim 1/\sqrt{N}$. The increase of the plateau with the storage load seems logical: the more patterns are stored, the less efficient the retrieval of any single one becomes.

\subsection{Dynamics visualization: initialization from a random state}
 We now aim at visualizing the evolution of the configurations when initializing the system from a random state. Again, the model is built on a $32 \times 32$ fully connected square lattice and each replica has its own saved patterns and initial state. The temperature is still set to $\beta=5.0$. We give the evolution of $m^\mu$ for all patterns for a given replica in \Cref{fig:all_mu_single_replica}.
\nolinenumbers
\begin{figure}[htbp]
\centering
\includegraphics[width=\columnwidth]{all_mu_single_run_N1024_beta=5.0_alpha0.030.pdf}
\caption{Evolution of $m^\mu$ for all patterns for the replica 16 of $\alpha=0.03$.}
\label{fig:all_mu_single_replica}
\end{figure}
\linenumbers
At the beginning, the state is roughly uncorrelated with  any of the patterns. As the system evolves, it starts correlating with some patterns, before it eventually aligns perfectly ($m^\mu=-1$) with one of them. When it does so, the overlaps with the other patterns collapse to 0, with a noise of standard deviation $1/32$. This configuration looks stable (as we can see small fluctuations around it), which was expected has it corresponds to a local minimum of energy.\\
We can also look at the evolution of the maximum overlap $m^*$ for all replicas at a given value of the storage load, as shown in \Cref{fig:mu_star_all_replica}.
\nolinenumbers
\begin{figure}[htbp]
\centering
\includegraphics[width=\columnwidth]{mu_star_all_runs_N1024_beta=5.0_alpha0.030.pdf}
\caption{Evolution of the maximum overlap of each replica for $\alpha=0.03$.}
\label{fig:mu_star_all_replica}
\end{figure}
\linenumbers
While in some case the state quickly falls into a stable minimum of energy corresponding to a specific pattern, some others seem to fluctuate around spurious states, for which the overlap is nor 1 or 0. While most configurations remain stuck in those spurious states (for the evolution time considered here), some others eventually escape those states to align with a specific pattern. The evolution obtain can be seen as the analogy of \Cref{fig:evolution_ising} in the case of the Hopfield network.\\
Lastly, we can plot the evolution of the mean of maximum overlaps across all replicas $\langle m^* \rangle$ depending on the storage load, see \Cref{fig:mean_mu_star}.
\nolinenumbers
\begin{figure}[htbp]
\centering
\includegraphics[width=\columnwidth]{mean_overlap_N1024_beta=5.0.pdf}
\caption{Evolution of $\langle m^* \rangle$ with the storage load. }
\label{fig:mean_mu_star}
\end{figure}
\linenumbers
As expected, for the lowest values of $\alpha$, the mean maximum overlap quickly increases to $\langle m^* \rangle \approx 1$ and remains stable. The plateau value decreases as the storage load increases: it is less and less likely for the system to end up in a completely pattern-aligned configuration from a random initialization when the number of spurious states increases.

\FloatBarrier
\subsection{Performance test}
We present in \Cref {fig:speedup_hopfield} the performance analysis of the code obtained with the implementation presented in \Cref{sec:base_implementation_hopfield}. Several values of $\alpha$ and lattice sizes have been tested for an evolution of $2^{10}$ sweeps each.
Subfigure (\subref{fig:sweep_time_vs_N_hopfield}) confirms that both implementations have complexities $\mathcal{O}\left(N \cdot N_\text{neighbors}\right)\sim \mathcal{O}\left(N^2\right)$ per sweep.

Compared to the two previous models, the speedup factor strikingly decreases close to 1 (and in some cases even below) in the case of the dense Hopfield network. The bitwise implementation is no longer more efficient than the reference one. In addition, the speedup factor appears to be mostly uncorrelated with the temperature and the storage capacity of the network.

The main explanation for this drop in efficiency of the bitwise implementation is the topology and the structure of the interactions in the case of the Hopfield model. Indeed, the Hopfield network is most commonly defined as a dense (fully connected) network, meaning that every neuron interacts with every other neuron in the network. This contrasts with the two previous models, in which each spin interacted with only four neighbors; in the Hopfield network, each neuron interacts with all $N-1$ remaining neurons. Thus, if the loop over neighbors at each spin update is not perfectly optimized in the bitwise implementation, the loss in performance can become considerable for dense networks, as observed here. The loss in performance is however much more discrete in the case of the Ising and spin glass models, for which we only sum on the 4 closest neighbors each time.

This hypothesis is confirmed by the decrease of the speedup factor with the lattice size, impacting the number of neighbors to sum over. Also, when testing the code on a sparse Hopfield model, the obtained speedup is very comparable to that of the other models. Conversely, when initialized on a dense network, the Ising model bitwise implementation also yields a significantly reduced speedup factor, confirming again our hypothesis. This is also reassuring, as it shows that the issue does not stem from the generalized use of \UnsignedInt in this new bitwise implementation, but rather a common lack of optimization in all models, that finds its sources in the coding of the base classes \Bits and \UnsignedInt.

The main reason for this lack of efficiency in the neighbors loop might the creation of temporary \texttt{Bits} objects over the summations, killing the code performance when done hundreds of thousands ($\sim N^2$) times per sweep.

Two other related causes can explain the lack of performances for dense networks. The first one is the data structure and the other one is the memory access. The structure has indeed initially been built for sparse networks, with the use of lists of neighbors instead of direct interaction matrices. This implementation is effective when the network is sparse but becomes inefficient in the case of a dense network: instead of storing a fully filled $N \times N$ triangular interaction matrix containing $N \left(N-1\right)/2$ values (taking advantage of the symmetry), we store twice as many values in $N$ lists of $N-1$ components, thus increasing the memory usage, and potentially reducing the performance if the memory overflows the available RAM.
\nolinenumbers
\begin{figure*}[p]
\centering
% ── Rangée 1 : a et b ────────────────────────────────────────────────
\begin{subfigure}[b]{0.49\textwidth}
  \centering
  \includegraphics[width=\columnwidth,height=0.17\textheight,keepaspectratio]{speedup_curves_hopfield.pdf}
  \caption{Speedup factor evolution.}
  \label{fig:speedup_curve_hopfield}
\end{subfigure}\hfill
\begin{subfigure}[b]{0.49\textwidth}
  \centering
  \includegraphics[width=\columnwidth,height=0.17\textheight,keepaspectratio]{time_normalized_hopfield.pdf}
  \caption{Normalized time evolution.}
  \label{fig:normalised_time_hopfield}
\end{subfigure}

\vspace{0.5em}
% ── Rangée 2 : c en pleine largeur ───────────────────────────────────
\begin{subfigure}[b]{\textwidth}
  \centering
  \includegraphics[width=\columnwidth,height=0.15\textheight,keepaspectratio]{time_total_vs_N_hopfield.pdf}
  \caption{Sweep time depending on the lattice size for all temperatures.}
  \label{fig:sweep_time_vs_N_hopfield}
\end{subfigure}

\vspace{0.5em}
% ── Rangée 3 : d à gauche ────────────────────────────────────────────
\begin{subfigure}[b]{0.49\textwidth}
  \centering
  \includegraphics[width=\columnwidth,height=0.17\textheight,keepaspectratio]{speedup_curves_by_alpha_hopfield.pdf}
  \caption{Speedup factor depending on the temperature for different $\alpha$ values.}
  \label{fig:speedup_by_alpha_hopfield}
\end{subfigure}\hfill
\hspace*{0.49\textwidth}

\vspace{0.5em}
% ── Rangée 4 : e en pleine largeur ───────────────────────────────────
\begin{subfigure}[b]{\textwidth}
  \centering
  \includegraphics[width=\columnwidth,height=0.15\textheight,keepaspectratio]{time_normalized_by_alpha_hopfield.pdf}
  \caption{Sweep time depending on the temperature for different $\alpha$ values.}
  \label{fig:sweep_time_vs_alpha_hopfield}
\end{subfigure}

\caption{Performance analysis of the bitwise and classical Hopfield Metropolis implementations. In Subfigures~(\subref{fig:speedup_curve_hopfield}), (\subref{fig:normalised_time_hopfield}) and (\subref{fig:sweep_time_vs_N_hopfield}), the data has been averaged over the different values of $\alpha$.}
\label{fig:speedup_hopfield}
\end{figure*}
\linenumbers
